\chapter{AI-Assisted Execution Workflow}
\ChapterLabel{AIImplementationWorkflow}
\index{AI agent!execution workflow}
\index{AI agent!parallel delegation}
\index{AI agent!independent review}

\section{Chapter Responsibility}
\label{Section:AIImplementationWorkflowResponsibility}
Organize AI-assisted work from planning through integration and human review.

\noindent\textbf{Expected result:} Coordinated execution with explicit checks and approval boundaries.

\section{Purpose and Scope}
\label{Section:AIImplementationPurpose}
Use this shared workflow for the AI skills selected in
\ChapterRef{AIAgentGuide}: project initialization, architecture documents,
requirements, user stories, analyses, code, tests, verification, validation,
and release-readiness recommendations. Appendix~C defines the task inputs,
context, steps, and acceptance criteria; this appendix defines how to plan,
assign, execute, integrate, check, and review the work with one or more agents.
Agent count does not change a task's acceptance criteria or human approval
boundaries. Roles describe responsibilities, not a required agent count.
Codex supplies the concrete examples; see
\hyperref[Section:AIOtherTools]{Other AI Tools} for adapting the workflow.

Follow the source-of-truth and human-decision rules in Appendix~C. Use approved
baselines when implementing code; document drafts and analyses may propose
changes but shall not claim human approval. Apply checks appropriate to the
artifact: document builds and reciprocal mappings for architecture content,
project builds and tests for code, controlled evidence for verification, and
stakeholder observations for validation. Project-specific commands and
completion criteria belong in the implementation and development-plan
chapters; controlled relationships remain in their authoritative registers.

Figure~\ref{Figure:AITaskExecutionActivity} connects
\ChapterRef{AIAgentGuide} to this workflow. Appendix~C supplies the selected
skill definitions; this workflow applies them with one or multiple agents and
checks the results against their Acceptance. The final check uses the
\hyperref[Table:AICompletionChecklist]{completion checklist}; completion does
not imply human approval. This overview summarizes the successful path.
Use Figure~\ref{Figure:AIImplementationSequence} and the
\hyperref[Section:AIImplementationVerification]{correction and escalation rules}
for failed checks and human decisions.

\section{Starting the Workflow in an AI Chat}
\label{Section:AIImplementationStart}
Open a Codex chat in the project workspace and replace the bracketed fields.
Starting delegated work involves three distinct responsibilities:
\begin{enumerate}
\item \textbf{Codex supplies the tools.} The application exposes tools to
create, message, wait for, and interrupt subagents. This appendix and its
prompts do not enable those tools. Current local Codex releases enable
subagents by default; if disabled, enable \texttt{agents.enabled} in the
Codex configuration before requesting delegated work.\footnote{Codex subagent
configuration and delegation guidance: \cite{CodexSubagents}.}
\item \textbf{The human authorizes delegation.} The request or applicable
project instructions permit subagents and set scope, budget, concurrency,
and decision boundaries. Tool availability alone is not authorization.
\item \textbf{The aiaCoordinator decides and invokes.} Confirm tool availability
and authorization, estimate the work, and choose one agent or independent
bounded assignments. If delegation is useful and permitted, invoke the
actual Codex subagent tools and wait for their results before integration.
\end{enumerate}
Use the starting request in the project chat; creating an API session or
running Python is not required. Naming AI agent roles in a prompt does not
launch subagents. Name authoritative sources and deliverables; the
aiaCoordinator reports actual agents used and project-specific commands.
See \hyperref[Section:AIAgentDelegation]{Choose and Invoke Agents} for assignment
and invocation requirements. The examples follow the
\hyperref[Section:AIPromptFormatting]{prompt formatting convention} in Appendix~C.

\begin{samepage}
\subsection{Start Any Selected AI Skill}
\label{Section:AIStartSelectedTaskPrompt}

\begingroup\ttfamily\raggedright
\noindent In [project workspace], complete [outcome] using AIS-[IDs] from Table C.2 in dsnAIAgentGuide.tex. Read their sources, Context, Steps, and Acceptance. Follow the shared execution workflow in dsnAIImplementationWorkflow.tex. Act as aiaCoordinator. For substantial work, first apply \hyperref[Skill:AIS-01]{\texttt{AIS-01}}: estimate effort and token ranges, assumptions, and dependencies. You are authorized to use Codex subagent tools for independent bounded tasks, after confirming their availability, within project instructions and a maximum of [N] concurrent agents including yourself. Explain the delegation decision and assign each agent its AI Skill, Context, Steps, output, Acceptance, dependencies, and file ownership. Use one agent when delegation adds little value; if subagents are unavailable, report that and work sequentially. Integrate outputs, run the applicable checks, and obtain independent review when warranted. Arrange stakeholder validation where needed. After two unsuccessful corrections of the same failure, pause dependent work and present evidence and options. Ask me to resolve material decisions or required cost approvals before dependent work. Report actual agents used, changes, commands, evidence, and gaps. Do not submit, approve a baseline, or authorize release for me.\par
\endgroup

\end{samepage}

\begin{samepage}
\subsection{Develop an Approved Change}

\begingroup\ttfamily\raggedright
\noindent In [project workspace], develop [approved change] against [approved architecture and requirements]. Read the AI agent guide and Appendix D. As aiaCoordinator, list components, criteria, checks, and decision gaps. Assign bounded tasks with file ownership. You are authorized to use Codex subagent tools for independent work, with at most [N] concurrent agents including yourself. Confirm tool availability, explain the delegation decision, and invoke parallel aiaDeveloper agents only when useful; otherwise work sequentially. Integrate, verify, and obtain independent review. Arrange stakeholder validation; record only observed results. After two failed corrections of one check, pause dependent work and show evidence and options to revise or stop. Ask me to decide material scope, architecture, safety, security, release, or baseline issues. Report files, commands, checks, trace links, validation evidence, assumptions, and gaps. Do not submit or approve the change for me.\par
\endgroup

\end{samepage}

The same sequence can prepare a requirements package before implementation.
For that use, aiaDeveloper agents produce bounded analyses and drafts, the aiaCoordinator
is the single writer for shared controlled registers, and checks include
traceability and document builds. Replace its placeholders with project facts.

\begin{samepage}
\subsection{Prepare a Requirements Package}

\begingroup\ttfamily\raggedright
\noindent In [project workspace], develop a reviewable requirements package for [approved scenario]. Read the approved architecture, stakeholder scenario, applicable domain profiles, existing registers, AI agent guide, and Appendix D. List missing source facts as questions; do not invent project facts. Act as aiaCoordinator. First show the sources, artifacts, acceptance criteria, and task boundaries. You are authorized to use Codex subagent tools, with at most four concurrent agents including yourself. Confirm tool availability and explain the delegation decision. If independent work warrants it, invoke up to three AI agents: aiaDeveloperA for stories and validation scenarios, aiaDeveloperB for functional, interface, and fault requirements, aiaDeveloperC for verification and test design. Keep one writer per shared file; if subagents are unavailable, work sequentially. Integrate the outputs, create reciprocal trace links, run document and build checks, and report actual results. Obtain an independent review of the integrated candidate. After two unsuccessful corrections of a failed check, pause dependent work and show me evidence and options. Ask me to decide material scope, architecture, safety, or baseline issues. Report changed files, trace gaps, assumptions, evidence, and open issues. Do not submit or approve a baseline for me.\par
\endgroup

\end{samepage}

\subsection{Other AI Tools}
\label{Section:AIOtherTools}
Codex is the concrete application used in the starting examples. The AI Skills,
Context, Steps, Acceptance, and responsibilities of aiaCoordinator,
aiaDeveloper, and aiaReviewer also apply when using other AI tools. The
application provides the execution capabilities; a model name alone does not
establish access to project files or tools for delegating work.

\begin{description}
\item[Claude Code.] Supports subagents with separate context and specified
instructions, tools, and permissions. Adapt the starting request to Claude
Code's available subagents and delegation tools; confirm the instructions
and sources each subagent receives.\footnote{Claude Code subagent documentation:
\cite{ClaudeCodeSubagents}.}
\item[Gemini.] Gemini CLI supports subagents exposed to the main agent as tools;
a delegated subagent returns its findings to the main agent. Confirm which
subagents and permissions are available in the installed application before
assigning work. Capabilities of Gemini CLI shall not be assumed for every
Gemini chat interface.\footnote{Gemini CLI subagent documentation:
\cite{GeminiCLISubagents}.}
\item[Grok.] Grok offers a multi-agent research feature that coordinates agents
for research and synthesizes their findings. Before using it for this workflow,
confirm whether the selected Grok application or API supports the required
project access, assignment boundaries, and checks. Research-agent collaboration
alone does not demonstrate that separate aiaDeveloper agents edited and verified
project artifacts.\footnote{Grok multi-agent research documentation:
\cite{GrokMultiAgent}.}
\end{description}

For any other application, replace ``Codex subagent tools'' in the example
request with its confirmed delegation capability. Keep the same AIS IDs,
Context, Steps, Acceptance, role names, human authorization, file ownership,
concurrency limits, verification, validation, and escalation rules. Identify
the application and version in the execution report and record actual agents
and checks used. If delegation is unavailable, report that limitation and
complete the applicable work with one agent. A prompt cannot enable missing
tools, and role names alone do not launch separate agents.

\section{Plan the Selected Tasks}
\label{Section:AIImplementationPlan}
Before changing artifacts, select the applicable AI Skills and identify the
requested outcome, approved constraints, relevant domain profiles, task
acceptance criteria, affected sources, and required checks. For code changes,
identify the approved requirements, architecture, interfaces, and tests. Use the AI Skills Register in
\ChapterRef{AIAgentGuide} to load the affected sources and reciprocal mappings.
Record any gap or conflict as a question with its implementation impact. An
authorized human resolves a material architecture, scope, safety, security, or
release decision before work that depends on that decision proceeds. An agent
may prepare independent analysis while the decision is pending; it shall not
silently choose an answer.

Apply AI skill \hyperref[Skill:AIS-01]{\texttt{AIS-01}} in
\hyperref[Table:AISkills]{the AI Skills Register} before substantial or
delegated work. Break the approved work into bounded tasks with inputs, expected outputs,
acceptance criteria, dependencies, and ownership of files or components. The
aiaCoordinator decides whether delegation will shorten the work or improve review
coverage enough to justify its coordination and token cost. Use one agent for
short tasks, ordered chains of dependent steps, or changes concentrated in the
same mutable files. Delegate in parallel only when tasks can make progress
independently. Give each aiaDeveloper a distinct edit boundary or an isolated
workspace; read-only investigations may share the same sources.

\section{Choose and Invoke Agents}
\label{Section:AIAgentDelegation}
\index{AI agent!agent-count decision}
No AIS inherently requires multiple agents. Use one agent for small edits,
a dependent chain of steps, or work concentrated in shared mutable registers.
Several AI Skills can be completed sequentially by one agent. Conversely, one AIS
can be partitioned among multiple aiaDeveloper agents when its outputs are independent.
Examples include separate architecture-view analyses, tests for independent
components, or separate requirements-family drafts. Dependent outputs wait
for their inputs; parallelism does not remove that dependency.

The human sets the outcome, authorization, budget, concurrency limit, and
project constraints. Within those limits, the aiaCoordinator decides whether to
delegate, after applying \hyperref[Skill:AIS-01]{\texttt{AIS-01}} when required. Explain the decision,
expected benefit, and coordination cost. Obtain human authorization when
project instructions require it, when delegation is not permitted, or when a
proposed plan exceeds agreed scope or cost. Do not interpret an estimate as
approval. If no delegation is authorized, perform the work with one agent.

To invoke multiple agents, the aiaCoordinator shall:
\begin{enumerate}
\item Confirm that the selected application exposes subagent tools and that the human request or
applicable project instructions authorize their use, as described in
\hyperref[Section:AIImplementationStart]{Starting the Workflow in an AI Chat}. Report its absence and use the
single-agent path when delegation is unavailable. For applications other than
Codex, apply \hyperref[Section:AIOtherTools]{Other AI Tools}.
\item Create one bounded assignment per aiaDeveloper with applicable AIS IDs,
authoritative source paths, linked Context, Steps, Acceptance, expected output, dependencies, and a distinct edit boundary. Reading
shared sources is allowed; controlled shared registers have one writer.
\item Launch assignments through the selected application's actual subagent tools, within the
total concurrency limit. Use the available messaging, waiting, and interruption
tools to coordinate their execution. A role name written in a
prompt is not evidence that a separate agent was launched.
\item Monitor returned results, resolve conflicts, and integrate ready outputs.
Record actual agents used, assignments, checks, and unresolved issues. Do not
claim parallel execution if the work was performed sequentially.
\end{enumerate}
For example, the aiaCoordinator may delegate story and validation-scenario
analysis under \hyperref[Skill:AIS-05]{\texttt{AIS-05}} and \hyperref[Skill:AIS-11]{\texttt{AIS-11}}, requirements drafting
under \hyperref[Skill:AIS-06]{\texttt{AIS-06}}, and test design under \hyperref[Skill:AIS-09]{\texttt{AIS-09}}. The aiaCoordinator
integrates their drafts into shared registers after dependencies are resolved.
These tasks may use several AI Skills; they do not create new skill IDs.

\section{Responsibilities and Handoffs}
\label{Section:AIAgentRoles}
The prefix \texttt{aia} means AI agent. These role identifiers follow
\hyperref[Rule:AITerminologyConsistency]{the terminology convention}.
\index{AI agent!aiaCoordinator}
\index{AI agent!aiaDeveloper}
\index{AI agent!aiaReviewer}

\begin{description}[leftmargin=3.1cm,style=nextline]
\item[aiaCoordinator] Confirms scope and decisions, partitions independent work,
assigns boundaries and acceptance criteria, reconciles results, integrates the
change, runs system-level checks, and reports evidence and open issues.
\item[aiaDeveloper] Produces the assigned analysis, document draft, code, tests, or evidence,
changes only the assigned area, preserves the approved architecture and project conventions, runs relevant local checks, and returns
a summary of changed files, results, assumptions, and unresolved issues.
\item[aiaReviewer] Examines the integrated result against requirements,
architecture, interfaces, tests, and evidence. The aiaReviewer reports concrete
findings and does not mark a decision or release approved on behalf of the
human owner.
\end{description}

Use the role names above consistently in instructions, figures, and reports.
The aiaDeveloper is the agent assigned to a bounded task, including document
authoring or analysis; the
number of aiaDeveloper agents does not change the role name. One agent may coordinate and develop a small change, then perform a separate
self-review pass. Report that pass as self-review, not independent review.
When independent review is required, use an uninvolved agent or human; if none
is available, report the review as pending. For a delegated change,
the aiaReviewer should be independent of its implementation and should review
only after integration produces a stable candidate. The aiaCoordinator resolves
conflicting aiaDeveloper results and repeats checks after corrections. Figure~\ref{Figure:AIImplementationSequence}
shows the optional parallel path and the later review handoff.

\section{Develop, Integrate, and Verify}
\label{Section:AIImplementationVerification}
Each aiaDeveloper keeps changes traceable to the approved task and supplies
local check results. The aiaCoordinator inspects and integrates all outputs,
checks shared interfaces and architecture constraints, updates the affected
documentation and reciprocal traceability, and runs the applicable complete build and checks: document compilation,
reference and mapping checks for documents; project builds and relevant tests
for code; evidence checks for verification and validation records. A failed check returns to the responsible task for
correction, followed by integration and verification again. Limit this cycle
to two unsuccessful correction attempts for the same failure. If the check
still fails, or the cause requires a decision outside the approved scope, the
aiaCoordinator reports the evidence and options to the human owner and pauses
dependent work until the owner decides whether to revise scope, continue with
a new bounded plan, or stop. A revised plan returns to task partitioning and
the decision about whether parallel delegation is useful; a stop ends the
workflow without review of an unverified candidate. Record actual
commands, results, controlled artifact or implementation version, and evidence locations;
do not infer a pass from a proposed test or an agent's assertion.

The review checks the integrated candidate for requirement
coverage, architecture conformance, unintended changes, missing tests, and
unsupported evidence. The aiaCoordinator addresses findings and obtains another
review when corrections affect the reviewed behavior. Human owners retain the
authority for the decisions named in \ChapterRef{AIAgentGuide}. A verified and reviewed candidate proceeds to stakeholder validation when
the task requires validation, with open issues reported. A document-only task
may finish after its acceptance criteria, checks, and required review are met;
it shall not claim implementation verification or stakeholder validation from
a successful document build.

\section{Validate with Stakeholders}
\label{Section:AIImplementationValidation}
When the selected task requires validation, prepare realistic scenarios mapped
to the approved user stories and use cases. Stakeholders evaluate the candidate
in its intended context against agreed success criteria. The aiaCoordinator may
prepare the scenarios and organize evidence, but records validation as passed
only from actual stakeholder observations and an authorized human disposition.
Update the authoritative validation register with the result, evidence
location, unmet needs, and requested changes. Passing implementation tests or
agent review alone does not establish validation.

If validation reveals a gap, the aiaCoordinator reports its impact and options to
the human owner. A material change to scope, requirements, or architecture
returns to planning and the delegation decision. A correction within approved
scope returns to bounded development; repeat affected integration,
verification, independent review, and validation before reporting completion.
The authorized human retains acceptance and release decisions.

\section{Project Tailoring Checklist}
\label{Section:AIImplementationTailoring}
When instantiating this template, record the following in the project-owned
implementation and development-plan chapters rather than in tool bootstrap
files:
\begin{itemize}
\item The AI tools permitted, how aiaDeveloper agents are launched, the maximum
useful concurrency, and whether isolated workspaces are required.
\item The project file or component ownership rules, integration owner, and
required code review or change-control process.
\item The exact build, static analysis, unit, integration, and system test
commands, including any hardware-in-the-loop or manual checks.
\item The required traceability, test evidence, implementation-version record,
and acceptance criteria for completion.
\item The validation scenarios, stakeholder participants, success criteria,
evidence locations, and authorized owner of the validation disposition.
\item The human decision owners and the points at which unresolved decisions
pause dependent implementation or release.
\end{itemize}

\begin{figure}[htbp]
  \centering
  \resizebox{\linewidth}{!}{\begin{tikzpicture}[
  x=1cm,y=1cm,
  participant/.style={draw=dsnBlue,fill=dsnLightBlue,rounded corners=2pt,
    minimum width=2.15cm,minimum height=0.55cm,align=center,font=\footnotesize},
  lifeline/.style={draw=dsnGray,densely dashed,line width=0.55pt},
  msg/.style={-{Latex[length=1.8mm]},draw=dsnBlue,line width=0.65pt},
  response/.style={-{Latex[length=1.8mm]},draw=dsnGray,dashed,line width=0.65pt},
  note/.style={font=\scriptsize,fill=white,inner sep=1pt},
  frame/.style={draw=dsnGray,rounded corners=2pt,line width=0.6pt}
]
  \node[participant] at (0,0) {Human owner};
  \node[participant] at (2.8,0) {aiaCoordinator};
  \node[participant] at (5.6,0) {aiaDeveloperA};
  \node[participant] at (8.4,0) {aiaDeveloperB};
  \node[participant] at (11.2,0) {aiaReviewer};
  \foreach \x in {0,2.8,5.6,8.4,11.2}
    \draw[lifeline] (\x,-0.35) -- (\x,-17.0);

  \draw[msg] (0,-1.0) -- node[note,above] {task, authorization, limits} (2.8,-1.0);
  \draw[msg] (2.8,-1.75) -- node[note,above] {material decision, if needed} (0,-1.75);
  \draw[response] (0,-2.35) -- node[note,above] {decision or revised scope} (2.8,-2.35);

  \node[draw=dsnBlue,fill=dsnLightBlue,rounded corners=2pt,
    font=\scriptsize,align=center,minimum width=3.4cm]
    (delegationDecision) at (2.8,-3.15) {Confirm Codex tools; select AI Skills\\estimate, partition, decide delegation};

  \draw[frame] (1.45,-3.85) rectangle (9.7,-7.45);
  \node[anchor=west,font=\scriptsize\bfseries,fill=white,inner sep=2pt]
    at (1.6,-3.85) {opt: available tools + authorized independent work};
  \draw[msg] (2.8,-4.5) -- node[note,above] {invoke: AI Skills, Context, task A} (5.6,-4.5);
  \draw[msg] (2.8,-5.05) -- node[note,above] {invoke: AI Skills, Context, task B} (8.4,-5.05);
  \draw[dsnGray,dotted] (1.65,-5.4) -- (9.5,-5.4);
  \node[anchor=west,font=\scriptsize\itshape,fill=white,inner sep=1pt]
    at (1.75,-5.4) {aiaDeveloper agents proceed concurrently in separate edit boundaries};
  \draw[response] (5.6,-6.2) -- node[note,above] {changes and local checks} (2.8,-6.2);
  \draw[response] (8.4,-6.85) -- node[note,above] {changes and local checks} (2.8,-6.85);

  \node[draw=dsnBlue,fill=dsnLightBlue,rounded corners=2pt,
    font=\scriptsize,align=center,minimum width=2.4cm]
    at (2.8,-8.1) {Integrate and run\\complete checks};
  \draw[frame] (1.45,-9.0) rectangle (5.2,-10.35);
  \node[anchor=west,font=\scriptsize\bfseries,fill=white,inner sep=2pt]
    at (1.6,-9.0) {loop: at most two corrections};
  \draw[msg] (2.8,-9.55) -- ++(1.35,0) -- ++(0,-0.4) --
    node[note,below] {correct and rerun checks} (2.8,-9.95);

  \draw[frame] (-1.15,-10.75) rectangle (4.5,-12.9);
  \node[anchor=west,font=\scriptsize\bfseries,fill=white,inner sep=2pt]
    at (-1.0,-10.75) {opt: still failing or decision needed};
  \draw[msg] (2.8,-11.3) -- node[note,above] {evidence and options} (0,-11.3);
  \draw[response] (0,-11.85) -- node[note,above] {revise plan or stop} (2.8,-11.85);
  \draw[msg] (0,-12.35) -- (-1.7,-12.35) -- (-1.7,-3.15) -- (delegationDecision.west);
  \node[font=\scriptsize\itshape,fill=white,inner sep=1pt,rotate=90]
    at (-1.7,-7.8) {revised plan};
  \node[anchor=west,font=\scriptsize\itshape,fill=white,inner sep=1pt]
    at (-0.9,-12.65) {stop: exit without review};

  \draw[msg] (2.8,-13.6) -- node[note,above] {if checks pass: candidate and evidence} (11.2,-13.6);
  \draw[response] (11.2,-14.25) -- node[note,above] {findings or review result} (2.8,-14.25);
  \draw[msg] (2.8,-14.9) -- node[note,above] {candidate and validation plan} (0,-14.9);
  \draw[response] (0,-15.65) -- node[note,above] {observations and disposition} (2.8,-15.65);
  \draw[msg] (2.8,-16.4) -- node[note,above] {validation record and open issues} (0,-16.4);
\end{tikzpicture}
}
  \caption[AI-assisted development and validation handoffs.]{AI-assisted development and validation handoffs. Codex supplies subagent tools; the human authorizes delegation; the aiaCoordinator decides and invokes. Parallel aiaDeveloper agents are used only for independent tasks; otherwise the aiaCoordinator develops tasks sequentially.}
  \label{Figure:AIImplementationSequence}
\end{figure}

\begin{figure}[htbp]
  \centering
  \resizebox{\linewidth}{!}{% Shared by the architecture manual and Week 03 slides.
\providecommand{\AIActivityRef}[2]{\hyperref[#1]{#2}}
\providecommand{\AIActivityCColor}{dsnBlue}
\providecommand{\AIActivityDColor}{dsnGreen}
\providecommand{\AIActivityFlowColor}{dsnGray}
\begin{tikzpicture}[
  action/.style={rounded corners=2pt,text width=4.15cm,minimum height=1.45cm,
    inner sep=5pt,align=center,font=\fontsize{10}{12}\selectfont},
  ctask/.style={action,draw=\AIActivityCColor,fill=\AIActivityCColor!6},
  execution/.style={action,draw=\AIActivityDColor,fill=\AIActivityDColor!6},
  flow/.style={-{Stealth[length=2mm]},draw=\AIActivityFlowColor,line width=.8pt}
]
\node[ctask] (select) at (0,0) {\textbf{Appendix C: skills}\\
  \AIActivityRef{Table:AISkills}{Select AI Skills; read Context}\\
  \AIActivityRef{Section:AIArchitectureWorkflow}{Read Steps and Acceptance}};
\node[execution] (plan) at (5.05,0) {\textbf{Appendix D: plan}\\
  \AIActivityRef{Section:AIImplementationPlan}{Plan the selected tasks}\\
  \AIActivityRef{Section:AIAgentDelegation}{Choose one or multiple agents}};
\node[execution] (develop) at (10.1,0) {\textbf{Appendix D: develop}\\
  \AIActivityRef{Section:AIImplementationVerification}{Develop assigned tasks}\\
  Use the selected skill definitions};
\node[execution] (check) at (10.1,-3.0) {\textbf{Appendix D: check}\\
  \AIActivityRef{Section:AIImplementationVerification}{Integrate, verify, and review}\\
  \AIActivityRef{Section:AIImplementationValidation}{Validate when required}};
\node[ctask] (accept) at (5.05,-3.0) {\textbf{Appendix C: Acceptance}\\
  \AIActivityRef{Section:AIArchitectureWorkflow}{Check task Acceptance}\\
  \AIActivityRef{Table:AICompletionChecklist}{Apply completion checklist}};
\node[ctask] (report) at (0,-3.0) {\textbf{Appendix C: report}\\
  \AIActivityRef{Table:AICompletionChecklist}{Report results and evidence}\\
  List open decisions};
\node[circle,fill=\AIActivityFlowColor,inner sep=2.5pt] (start) at (-3,0) {};
\node[circle,draw=\AIActivityFlowColor,inner sep=3.5pt] (end) at (0,-4.5) {};
\node[circle,fill=\AIActivityFlowColor,inner sep=2pt] at (end.center) {};
\draw[flow] (start) -- (select);
\draw[flow] (select) -- (plan);
\draw[flow] (plan) -- (develop);
\draw[flow] (develop) -- (check);
\draw[flow] (check) -- (accept);
\draw[flow] (accept) -- (report);
\draw[flow] (report) -- (end);
\end{tikzpicture}
}
  \caption{Activity connecting Appendix C skill definitions and Acceptance to Appendix D execution. The agent-count decision belongs to Appendix D for every selected task.}
  \label{Figure:AITaskExecutionActivity}
\end{figure}
